\section{Exact rational realization and bit complexity}
\label{sec:arithmetic}

The construction in the preceding sections counts local operations.  For an
explicit polynomial input, those operations can be performed exactly with
rational arithmetic.  The endpoint rule used below is essential: the midpoint
of a leaf is used to form its model, whereas a backtracked local minimizer is
chosen at an endpoint.  Consequently, repeated regridding does not repeatedly
multiply unrelated coordinate denominators.

We assume that every original interval and the initial center have rational
endpoints, and that the bag functions are given explicitly as sparse rational
polynomials
\begin{equation}
 a_t(v)=\sum_{\alpha\in\mathcal A_t}c_{t,\alpha}v^\alpha,
 \qquad c_{t,\alpha}\in\Q,
 \qquad \abs{\alpha}:=\sum_i\alpha_i\le D.
 \label{eq:arithmetic-polynomial}
\end{equation}
The total degree bound $D$ is fixed.  Let $S=\sum_t\abs{\mathcal A_t}$, and let
$L$ be the binary input length, including the tree, variable incidences,
polynomial coefficients and exponents, original intervals, initial center,
and the rational target accuracy $\eps>0$.  A supplied rational smoothness
bound is also included in $L$.  Rational numbers are represented by signed
integer numerators and positive integer denominators.  Fixed coordinates have
already been removed, as in the model section; the zero-dimensional case
requires only one exact evaluation.  Thus we take $p\ge1$ below.
Let $C_{\mathrm{den}}$ be the least common multiple of all polynomial
coefficient denominators; it is $1$ when there are no nonzero coefficients.

Bit work uses a random-access store of binary data, charging integer
arithmetic, copying, comparisons, and address handling by their bit cost.
The polynomial-time conclusions also hold under a standard Turing-machine
simulation, with polynomial overhead.

\subsection{A certified affine model}

Put
\[
 R=\max\{1,\abs{\text{an original endpoint}}\},
 \qquad H=\max\{1,\abs{c_{t,\alpha}}:t,\alpha\}.
\]
Both quantities are rational and can be computed by exact comparisons.  They
bound coordinate and coefficient magnitudes, respectively.

\begin{proposition}[A computable smoothness bound]
\label{prop:arithmetic-hessian}
For the polynomial input in \cref{eq:arithmetic-polynomial}, define
\[
 M_0=\max_t\sum_{\substack{\alpha\in\mathcal A_t\\\abs{\alpha}\ge2}}
 \abs{c_{t,\alpha}}\abs{\alpha}(\abs{\alpha}-1)
 R^{\abs{\alpha}-2}.
\]
Then $M_0$ is a nonnegative rational number, and
$\norm{\nabla^2a_t(v)}_2\le M_0$ on every original bag box.  For fixed $D$,
$M_0$ has $O_D(L)$ bits and is computable in polynomial bit time.  One may use
$M=M_0$ or any supplied rational $M\ge M_0$; the latter inequality is
checked exactly.
\end{proposition}

\begin{proof}
For a monomial of degree $d\ge2$, differentiation twice and summation over all
ordered coordinate pairs gives
\[
 \sum_{i,j}\abs{\partial_i\partial_j(c v^\alpha)}
 \le \abs{c}R^{d-2}
 \left(\sum_{i\ne j}\alpha_i\alpha_j+
       \sum_i\alpha_i(\alpha_i-1)\right)
 =\abs{c}d(d-1)R^{d-2}.
\]
The spectral norm of a matrix is at most the sum of the absolute values of
its entries.  Applying the triangle inequality to the Hessian proves the
claimed bound.  Terms of degree zero or one have zero Hessian.

The binary length of $C_{\mathrm{den}}$ is at most the sum of the coefficient
denominator lengths.  All terms in the displayed bound can be put over a
denominator dividing $C_{\mathrm{den}}$ times a fixed
power of the denominator of $R$.  Since $D$ is fixed, this denominator and
the summed numerator have $O_D(L)$ bits.  Exact addition, comparison, and the
computation of the required least common multiples therefore take polynomial
bit time.  If desired, a positive $M_0$ can be rounded upward to a power of
two, with a factor of less than two increase; $M_0=0$ is left unchanged.
\end{proof}

For a rational leaf $B$, let $m_B$ be its coordinatewise midpoint.  Use the
affine model
\begin{equation}
 \ell_{t,B}(v)=a_t(m_B)+\inner{\nabla a_t(m_B)}{v-m_B}
                  -\frac{Mp}{8}\width(B)^2.
 \label{eq:arithmetic-taylor}
\end{equation}
This notation $\ell_{t,B}$ denotes a bag model, and is distinct from the
separator affine function $l_{t,D}$ of the certificate definition.
Taylor's formula and the Hessian bound give
\[
 \abs{a_t(v)-a_t(m_B)-\inner{\nabla a_t(m_B)}{v-m_B}}
 \le \frac M2\norm{v-m_B}^2
 \le\frac{Mp}{8}\width(B)^2.
\]
Consequently
\[
 0\le a_t(v)-\ell_{t,B}(v)
       \le\frac{Mp}{4}\width(B)^2,
 \qquad v\in B.
\]
Thus \cref{eq:model-contract} holds with $A=Mp$.  The model is affine and
hence convex.  All its coefficients are obtained by exact rational
evaluation of the explicit polynomial and its derivatives.  A smaller
externally supplied smoothness bound can also be used if it has a separate
sound certificate.  The size and verification cost of that additional
certificate must then be included; the bounds below use $M\ge M_0$.

\begin{proposition}[Exact local minimization]
\label{prop:arithmetic-endpoints}
With the models in \cref{eq:arithmetic-taylor}, every local minimization in
the dynamic program of \cref{prop:model-dp} is the minimization of an affine
function on a rational box.  It can be solved exactly without a general
optimization oracle.  A deterministic minimizer chooses the lower endpoint
of each coordinate interval when its affine coefficient is nonnegative,
and the upper endpoint when its coefficient is negative.
\end{proposition}

\begin{proof}
The common separator slope $\lambda_t(c)$ is fixed before the stage.  For a
child $u$ and parent leaf $B$, the prescribed minorant has the form
\[
 \psi_{u,B}(s)=\inner{\lambda_u(c)}s+b_{u,B},
 \qquad
 b_{u,B}=\min\{\beta_{u,E}: E\text{ touches }B_{S_u}\}.
\]
It is affine, and its intercept is at most the intercept of every touching
child message.  The local objective consists of the bag model, these child
minorants, and the subtraction of the own separator affine slope.  It is
therefore affine.  Its domain is $B$ intersected with the cylinder over an
own separator cell $D$.  An intersection of coordinate boxes is a coordinate
box; its endpoints are maxima of lower endpoints and minima of upper
endpoints.  Empty intersections are omitted.

An affine function separates into the sum of one-coordinate affine
functions.  Each is minimized by the stated endpoint rule, including the
choice of the lower endpoint when the coefficient is zero.  Exact rational
sign comparisons determine all choices.  Taking minima over the finite
leaf and cell lists then gives the exact message intercepts and the root
lower bound.  Ties in those finite minima may be resolved by list order.
\end{proof}

This proposition also supplies exact local certificate inequalities: the
minimum on each nonempty leaf--cell intersection is evaluated at the stated
corner.  An arbitrary rational minimizer in a zero-slope direction would
still be mathematically valid, but would provide no bound on its denominator.
We use the deterministic endpoint rule throughout the arithmetic realization.

\subsection{Denominators and finite message sizes}

Let $Q$ be the least common multiple of the denominators of the original
endpoints, initial center, and $s_0$.  The last inclusion is redundant when
$s_0$ is the largest difference of two original endpoints, but makes the
invariant explicit.  Write $M=n_M/q_M$ with $q_M>0$, taking $q_M=1$ when
$M=0$.  Use a fixed integer grading exponent $\mu$ and
$\theta=2^{-\mu}$ with $\mu\ge1$ as in \cref{lem:regridding-shell}.

\begin{lemma}[Coordinate denominators]
\label{lem:arithmetic-denominators}
At stage $j$, every leaf endpoint, own separator cell endpoint, endpoint of
a local intersection, and reconstructed center coordinate has denominator
dividing $Q2^{j+\mu}$.  Every leaf midpoint has denominator dividing
$W_j=Q2^{j+\mu+1}$.  Their numerator and denominator lengths are
$O(L+j+\mu)$.
\end{lemma}

\begin{proof}
The initial center and original endpoints have denominators dividing $Q$.
For the induction step, the central partition endpoints are the current
center coordinate and its shifts by $h_j=s_0 2^{-j}$.  A shell at level
$\ell\le j$ starts at a shift by $-2^\ell h_j$ and has step
$2^{-\mu}2^{\ell-1}h_j$.  Hence every untruncated endpoint has the form
\[
 c_i+s_0 t\,2^{-(j+\mu)},\qquad t\in\mathbb Z.
\]
If $c_i=x_i^{j-1}$ has denominator dividing $Q2^{j-1+\mu}$, this new
endpoint has denominator dividing $Q2^{j+\mu}$.  Clipping to an original
interval selects an original endpoint when clipping is active.  Projections
and local intersections likewise select endpoints already present in the
stage.  They preserve the same divisibility.

By \cref{prop:arithmetic-endpoints}, each backtracked local minimizer consists
of such endpoints.  Reconstruction takes a coordinate from its top
occurrence bag, so it also selects one of these endpoints and closes the
induction.  A midpoint introduces at most one additional factor of two.
Midpoints are used only for model evaluation and are not substituted for
the reconstructed center.  Finally, all these points lie in the original
box and have magnitude at most $R$.  Since $\log_2 Q$ and $\log_2 R$ are
$O(L)$, the numerator bounds follow as well as the denominator bounds.
The initial stage obeys the same stated bounds, even though it need not
introduce the factor $2^\mu$.
\end{proof}

Define $d=\max\{D,2\}$ and the stage denominator
\begin{equation}
 Z_j=8q_M C_{\mathrm{den}} W_j^d.
 \label{eq:arithmetic-denominator}
\end{equation}
Here and below a denominator ``divides $Z_j$'' concerns the reduced rational
number; the implementation can store its integer numerator over the common
denominator $Z_j$ instead of repeatedly forming products of denominators.

\begin{lemma}[Common denominator and magnitude bound]
\label{lem:arithmetic-size}
Every model coefficient, common separator slope, finite message intercept,
child-minorant intercept, local minimum, root lower bound, and evaluated
incumbent objective at stage $j$ has denominator dividing $Z_j$.  Put
\[
 \mathcal H=(S+1)HR^D,
 \qquad
 V=(1+2D+3NpkD)\mathcal H+NMpR^2.
\]
The absolute values of finite message intercepts and objective values are
bounded by $V$, and all scalar quantities used in their calculation have
\begin{equation}
 b_j=O_D(L+j+\mu)
 \label{eq:arithmetic-scalar-bits}
\end{equation}
bits.  In particular, summation through a subtree does not cause exponential
growth of numerator or denominator lengths.
\end{lemma}

\begin{proof}
A midpoint coordinate has denominator dividing $W_j$.  Evaluation of a
degree-$D$ polynomial has denominator dividing $C_{\mathrm{den}}W_j^D$.
A derivative has denominator dividing $C_{\mathrm{den}}W_j^{D-1}$ when
$D\ge1$, and is zero when $D=0$.
Multiplication of a derivative by a midpoint or an endpoint returns a
denominator dividing $C_{\mathrm{den}}W_j^D$.  The model's constant correction has
denominator dividing $8q_M(Q2^{j+\mu})^2$.  Each of these denominators
divides $Z_j$.  The same argument applies to derivative evaluations at the
stage center used in $\lambda_t(c)$.

All remaining dynamic-programming operations are addition, subtraction,
endpoint evaluation, and selection of a finite minimum.  Products of a
separator slope with an endpoint have the same polynomial-evaluation
denominator just considered.  Induction over the tree therefore proves
divisibility by the single $Z_j$ for all finite entries.  An empty-domain
value is a symbolic $+\infty$ and requires no rational representation.

For completeness, the magnitude bound follows from the unfolded
configuration, rather than from repeated loose bounds on child messages.
The sum of absolute polynomial values at bag midpoints is at most
$\mathcal H$, and the sum of absolute gradient--displacement terms is at
most $D\mathcal H$.  A bag correction has magnitude at most $MpR^2/2$.
For any coordinate, a common separator slope sums derivatives from at most
$k$ bags, so its magnitude is at most $kD\mathcal H/R$.  A difference
between two coordinate copies has magnitude at most $2R$.  Thus the
separator terms in \cref{eq:model-config} have total absolute value at most
$2NpkD\mathcal H$.  The analogous expansion for a message intercept has
one additional, unpaired own-separator term, of magnitude at most
$pkD\mathcal H$.  The displayed $V$ bounds these quantities and the finite
minima over them.  It also bounds the individual polynomial and affine
quantities needed to form these entries.

Now $\log_2 Z_j=O_D(L+j+\mu)$.  The logarithm of $1+V$ is $O_D(L)$:
$N,p,k,S$ are bounded by the explicit input size, while $H,R$, and the
supplied or computed $M$ have $O_D(L)$ bits.  Representing a scalar over
$Z_j$ therefore needs $O_D(L+j+\mu)$ bits.  A sum of contributions from
$N$ bags adds only the logarithm of the corresponding magnitude factor;
it does not multiply their common denominators.
\end{proof}

\subsection{Bit work and certificate representation}

Put $m=4/\theta$, the scalar shell count used in
\cref{thm:regridding-main,prop:regridding-incidence}.  We retain its dependence
on the numerical grading $\theta$ explicitly.  Stage $j$ has at most
$m^p(j+1)$ leaves per bag.  The incidence construction of
\cref{prop:regridding-incidence} gives at most
$N3^q m^p(j+1)^2$ nonempty leaf--own-cell pairs, where $q\le p$.
We count polynomial terms as part of the input, rather than treating their
evaluation as one oracle call.  The work bound below uses the coarser
factor $3^p$ in place of $3^q$.

\begin{theorem}[Bit realization of a sufficiently graded run]
\label{thm:arithmetic-bit}
Use the affine models in \cref{eq:arithmetic-taylor} and the deterministic
endpoint rule.  Suppose the quadratic-growth assumption holds with some
$g>0$, and choose a dyadic grading satisfying
\cref{eq:regridding-theta}.  In the constants of
\cref{eq:regridding-constants}, set
\[
 \begin{aligned}
 A&=Mp, & C_0&=k(k-1)p,
 & B_0&=\frac{MC_0}{2}+\frac{Mp}{4},\\
 \eta&=\frac g{20},
 & C&=12B_0+\frac{6kC_0M^2}{\eta},
 & B&=\max\left\{p,\frac{2C}{g}\right\}.
 \end{aligned}
\]
Put
\[
 J=\max\left\{0,
 \ceil{\log_2\left(s_0\sqrt{\frac{gBN}{\eps}}\right)}\right\}.
\]
The certificate test $F(x)-\LB\le\eps$ succeeds by stage $J$.  There is an
exact rational implementation whose total bit work through that stage is
\begin{equation}
 \operatorname{poly}_D(L)+
 O_D\!\left(
 p(N+S)3^p m^p(J+1)^3\,\mathfrak b_J^2
 \right),
 \qquad
 \mathfrak b_J=O_D\!\left(L+J+\mu+p\log(m+1)+\log(J+2)\right).
 \label{eq:arithmetic-work}
\end{equation}
The logarithmic terms in $\mathfrak b_J$ also allow for table identifiers
and counters.  Scalar rational values alone obey the smaller bound
\cref{eq:arithmetic-scalar-bits}.
\end{theorem}

\begin{proof}
The analytic certificate and error bounds apply because
$\ell_{t,B}$ satisfies the model contract with $A=Mp$, and all local
minima are exact.  By \cref{thm:regridding-main,eq:regridding-distance}, stage
$j$ has certificate gap at most $gBNh_j^2$.  The chosen $J$ makes this at
most $\eps$.  Actual incumbent values and the certificate test are evaluated
exactly, so no additional rounding allowance is required.

Polynomial and gradient evaluation at all leaf midpoints uses
$O_D(Sm^p(j+1)+pNm^p(j+1))$ scalar operations: each explicit monomial has
at most $D$ factors counted with multiplicity.  Common slopes at the stage
center are computed from the sparse derivatives and the tree.  The local
affine corner evaluation and the coordinate comparisons for each nonempty
pair require $O_D(p)$ further operations.  Touching-cell minima and the
remaining message operations use the incidence lists.  Thus a conservative
stage bound is
\[
 O_D\bigl(p(N+S)3^p m^p(j+1)^2\bigr).
\]
This includes the explicit term count and does not assume an uncharged
polynomial-evaluation oracle.

For arithmetic, write stage coordinates over $Q2^{j+\mu}$, midpoints over
$W_j$, and scalar costs over $Z_j$.  Monomial numerators are obtained by
integer multiplication and exact scaling to these denominators.  No
coordinate division or optimization-dependent denominator is introduced.
The common-denominator representation also makes a message addition an
integer addition and a message comparison an integer comparison.  Exact
integer division, when used for scaling, and multiplication can be performed
in $O(\mathfrak b_J^2)$ bit operations by elementary algorithms.  The
coordinate tuples and list identifiers have polynomially bounded length
included in $\mathfrak b_J$; their coordinatewise comparisons are charged
to the factor $p$ in the stage bound. Table addresses fit in
$\mathfrak b_J$ bits. Balanced dictionaries keyed by grid-index tuples
give $O(p\mathfrak b_J^2)$ bit work per incidence lookup, also covered
by this bound. Preprocessing the input denominators
and the Hessian bound takes $\operatorname{poly}_D(L)$ bit operations.
Finally, summing $(j+1)^2$ for $0\le j\le J$ gives
$O((J+1)^3)$ and proves \cref{eq:arithmetic-work}.
\end{proof}

The numerical parameters in this theorem are not hidden inside a polynomial
in the input length.  In particular, $m$ depends on the sufficient grading,
which depends on $M/g$, $p$, and $k$, and the number of local records has the
factor $m^p$.  The constants in $J$ likewise retain their stated dependence
on $g,M,N,p,k,s_0$, and $\eps$.  For example,
$F(x_1,x_2)=x_1^2+2^{-b}x_2^2$ on $[-1,1]^2$ has a rational input of
$O(b)$ bits, Hessian bound $M=2$, and quadratic-growth constant $g=2^{-b}$.
The ratio $M/g$ is exponential in $b$.  The theorem is a bit realization
of the conditional numerical work bound; it does not assert polynomial time
in $L$ when that conditioning or the bag dimension is unrestricted.

\begin{proposition}[A compact exact certificate]
\label{prop:arithmetic-certificate}
The certificate of a sufficiently graded stage $J$, together with its
incumbent, has a rational serialization of
\[
 O_D\bigl(L+pNm^p(J+1)b_J\bigr)
\]
bits.  Given this serialization and the original input, all local certificate
inequalities and the final gap can be verified by exact rational operations
with the same per-stage type of bound as in
\cref{eq:arithmetic-work}.
\end{proposition}

\begin{proof}
Store the stage index $J$, center and grading, leaf endpoint tuples, own
separator cell endpoint tuples and intercepts, the shared slope
$\lambda_t(c)$ for each nonroot bag, the intercept $b_{u,B}$ for each
parent-leaf--child edge, the root lower bound, and the retained incumbent.
The affine bag model is reconstructed from the midpoint formula, so
its coefficients need not be repeated in the serialization.  Child
minorants use their already stored shared slopes.  Likewise, local corner
minima and proofs for every leaf--cell pair need not be stored: a verifier
recomputes them.

There are at most $Nm^p(J+1)$ leaves and no more own separator cells under
the construction of \cref{lem:regridding-shell}.  Moreover,
\[
 \sum_t\abs{\operatorname{ch}(t)}\abs{\mathcal L_t}
 \le (N-1)m^p(J+1).
\]
Thus the endpoint, slope, intercept, and incumbent data comprise
$O(pNm^p(J+1))$ scalar entries.  Each has $O_D(b_J)$ bits by
\cref{lem:arithmetic-denominators,lem:arithmetic-size}.  List references may
be encoded as coordinate-index tuples and fit in the same total bound.
Including the original input once gives the stated serialization size.

A verifier first computes $M_0$, checks $M\ge M_0$ and that the stored
stage center belongs to $X$, then checks the generated partitions,
midpoint models, and shared slopes. For each child
and parent leaf it checks the minorant
intercept against every touching own cell.  It then evaluates the local
affine minimum at the deterministic corner of every nonempty own-cell
intersection and checks the corresponding certificate inequality.  The
incidence bound controls the number of these checks.  Finally, it checks
that every incumbent coordinate belongs to its original interval, and
evaluates the incumbent objective, the root lower bound, and their
difference exactly.  Validity then follows from
\cref{prop:model-validity}; the verifier
does not need a certificate of the quadratic-growth constant $g$.
\end{proof}

\subsection{An unknown growth constant}

The positive number $g$ may be a promise used in the convergence analysis
rather than part of the rational input.  The stopping test does not use it:
every correctly formed certificate gives $\LB\le f^*$, and an actual
incumbent gives $F(x)\ge f^*$.  The inequality $F(x)-\LB\le\eps$ therefore
certifies the required accuracy for every grading, including a grading too
coarse for the convergence proof.  The smoothness bound is different: it
is needed for model validity itself and is computed or certified before
any run is accepted.

\begin{proposition}[Operation-budget search without a supplied $g$]
\label{prop:arithmetic-search}
For the rational polynomial input, suppose the quadratic-growth assumption
holds for some unknown $g>0$.  Run the dyadic grading candidates by the
following deterministic schedule.  In round $r=1,2,\ldots$, start each
candidate $1\le\mu\le r$ from the original center, and give it at most $2^r$
counted bit operations, including final serialization; accept the first
candidate that passes the exact certificate-gap test and finishes writing
its certificate within its budget.

For any sufficient candidate $\mu_*$, let $T_*$ be a bound of the form
\cref{eq:arithmetic-work} for its run through the corresponding $J$, enlarged
if necessary so that $T_*\ge\max\{2,2^{\mu_*}\}$.  The search terminates,
returns a certified $\eps$-accurate incumbent, and uses
$O(T_*\log T_*)$ total bit work.  It does not need a stage cap or an estimate
of $g$.
\end{proposition}

\begin{proof}
Each candidate uses the already valid smoothness bound and exact local
arithmetic.  Therefore any successful gap test is sound, independently of
whether that candidate satisfies the sufficient grading inequalities.
The existence of a positive $g$ guarantees a finite sufficient integer
$\mu_*$ because all inequalities in \cref{eq:regridding-theta} hold when
$\theta$ is sufficiently small.  Let
$R_* = \ceil{\log_2 T_*}$.  The enlargement of $T_*$ ensures
$\mu_*\le R_*$.  Round $R_*$ includes that candidate and supplies at least
$T_*$ bit operations, so it completes its successful run unless an earlier
candidate has already succeeded.

The sufficient-run bound covers serialization: its cost is at most the
output length in \cref{prop:arithmetic-certificate}, which is bounded by
\cref{eq:arithmetic-work}. The schedule restarts candidates and hence
requires no stored suspended states. Charging the bounded runs to their
operation budgets gives
\[
 \sum_{r=1}^{R_*}r2^r
 =(R_*-1)2^{R_*+1}+2
 \le 4T_*\ceil{\log_2 T_*}.
\]
These limits are elementary-step limits for the fixed exact-arithmetic
algorithm.  A binary countdown enforces a trial budget: over $2^r$
decrements, its total digit changes and head movements are $O(2^r)$,
since the successive digit positions are changed at most
$2^r,2^{r-1},\ldots,1$ times, respectively.  A maintained end marker detects
exhaustion without scanning the full counter at every step.  Initializing
the counter costs $O(r)$ per trial, and round enumeration is likewise
covered by the displayed sum.  Hence scheduling adds only a constant
factor, rather than another logarithmic factor.  Stopping an unsuccessful trial
does not assert any mathematical conclusion about it.  The input
preprocessing is done once and is already included in $T_*$.
\end{proof}

The logarithmic factor is an upper bound for this particular simple search,
not a lower bound on all ways of handling unknown parameters.  The first
successful candidate may use a different grading and a different stage
from $\mu_*,J$.  Its accuracy follows from its own certificate test, and
the preceding proposition bounds its work by the search budget.  The compact
size bound for a prescribed sufficiently graded stage should therefore not
be transferred to that first successful output without its actual grading
and stage.  A serialization of the actual output is always bounded by the
operations used to produce it, in addition to the input it references.

Finally, guessing a smoothness bound and applying the same stopping rule is
not justified: an insufficient smoothness bound can make a purported lower
bound invalid.  The explicit polynomial Hessian bound removes this issue.
The search varies only the grading; every candidate uses valid models from
the outset.
